% =============================================================================
\section{Dimension 2: The Router Layer}
\label{sec:router}
% =============================================================================

The router sits between the client and the inference pool.  Our
publications cover three adaptation regimes: static rules, online
bandit adaptation, and RL-based selection.

% -----------------------------------------------------------------------------
\subsection{Static Semantic Routing}
\label{sec:router:static}
% -----------------------------------------------------------------------------

The \vsr{} position paper~\cite{vllmsr2026} established the
signal-driven architecture: composing thirteen signal types into
deployment-specific policies across vLLM, OpenAI, Anthropic, Azure,
Bedrock, and Gemini backends.  The embedding similarity signal is
powered by mmBERT-Embed~\cite{mmberted2025}, whose 2D Matryoshka
design lets operators dial latency vs.\ quality per deployment
(2.56\,ms at layer-6 for real-time routing, 11\,ms at layer-22 for
maximum recall).  ProbPol~\cite{probpol2026} formalized conflict
detection for these signals.

Token-budget pool routing~\cite{poolrouting2026} implements a
specialized static rule: dispatching to short or long pools based on
estimated total token budget, achieving 17--39\% GPU reduction.
OATS~\cite{oats2026} addresses tool selection with zero-cost
embedding refinement, improving NDCG@5 from 0.869 to 0.940.

Separately, WhenToReason~\cite{wang2025reason} classifies whether a
query warrants reasoning-mode invocation, avoiding unnecessary
chain-of-thought tokens on simple queries.  The Factcheck
Classifier~\cite{factcheck2026} adds another pre-route signal:
queries classified as fact-seeking can be steered toward models with
lower hallucination rates or flagged for post-response validation
via HaluGate~\cite{halugate2025}.

For multi-turn sessions, the \vsr{} project~\cite{vllmsr2026}
provides two mechanisms beyond single-request classification.
First, \emph{context memory injection}: the router accumulates
contextual signals from prior turns (topic, domain, user profile,
difficulty trajectory) and injects this context into the routing
decision for subsequent turns, giving the inference endpoint richer context that improves response
accuracy.  This is similar in spirit to the conversation-state
management in multi-turn APIs like OpenAI's Responses API.  Second, \emph{contrastive embedding classification
for multi-turn safety}: rather than classifying each turn in
isolation, the router uses few-shot contrastive embeddings to detect
jailbreak attempts that spread adversarial content across multiple
turns, which single-turn classifiers cannot detect.

Content safety and privacy form another pre-route gate.  The
MLCommons Safety Level-1 classifier~\cite{safetyl1_2026} performs a
binary safe/unsafe check on every inbound request; flagged requests
are passed to the Level-2 hazard classifier~\cite{safetyl2_2026},
which assigns one of nine MLCommons categories (violent crimes,
non-violent crimes, sex crimes, weapons/CBRNE, self-harm, hate,
specialized advice, privacy, misinformation).  This hierarchical
design keeps latency negligible for safe traffic (only unsafe
requests pay the second-stage cost) and enables per-category routing
policies: privacy-violating prompts can be blocked at the gateway,
misinformation queries routed to fact-grounded models, and
specialized-advice requests logged for compliance auditing.  Both
classifiers use mmBERT LoRA adapters, fitting the project's
low-resource design constraint.

FastRouter~\cite{fastrouter2026} ensures all these routing
decisions---single-turn and multi-turn---happen within the latency
budget: 50\,ms end-to-end with an 800\,MB GPU footprint, small
enough to co-locate with serving.


% -----------------------------------------------------------------------------
\subsection{Online Adaptive Routing for Chat}
\label{sec:router:online}
% -----------------------------------------------------------------------------

Chat workloads offer a natural feedback signal: the user's next message indicates whether the previous response was
satisfactory.  Our
mmBERT-32K Feedback Detector~\cite{feedbackdet2026} classifies each
user turn into four categories (satisfied, needs clarification, wrong
answer, wants different approach) at 98.8\% accuracy with
$<$1\,ms overhead.  When the detector flags dissatisfaction, the
router can re-route subsequent turns to a stronger model or a
different pool---no explicit reward model is needed because the
user's own reaction serves as the reward.  Combined with the
router's context memory injection (\Cref{sec:router:static}), the
adaptation becomes session-aware: the router knows not just that the
user is dissatisfied, but which topic and difficulty level triggered
the dissatisfaction, enabling targeted re-routing rather than a
blanket model upgrade.

Beyond our project, RouteLLM~\cite{ong2025routellm} learns routers
from human preference data (95\% GPT-4 quality at 26\% cost);
MixLLM~\cite{lu2025mixllm} uses contextual bandits with query tags
(97.25\% of GPT-4 quality at 24.18\% cost);
SageServe~\cite{jia2025sageserve} integrates traffic forecasting
with ILP-based routing (25\% GPU-hour savings at Microsoft scale).


% -----------------------------------------------------------------------------
\subsection{RL-Based Routing for Agents}
\label{sec:router:rl}
% -----------------------------------------------------------------------------

Agent workloads lack the per-turn user feedback that chat enjoys:
the ``user'' is another LLM or an orchestrator, and task success is
only observable at the end of a multi-step session.  This delayed,
sparse reward makes RL a better fit than bandit methods.
Router-R1~\cite{routerr1_2025} demonstrates the approach: it
instantiates the router itself as an LLM, interleaving ``think''
actions (internal reasoning) with ``route'' actions (model
invocation), and trains end-to-end via PPO with a composite
reward---format, outcome, and cost---that lets the router learn
performance--cost trade-offs without hand-crafted heuristics.
R2-Router~\cite{r2router2026} jointly selects model and output
length budget at 4--5$\times$ lower cost.
M-CMAB~\cite{mcmab2026} uses online RL for multi-modal scheduling
under heterogeneous budgets.

Our AVR system~\cite{avr2026} applies cascading to VLM agent tasks:
Tier~1 multimodal classifier ($\sim$10\,ms), Tier~2 small-VLM
confidence probing with token-filtered scoring, Tier~3 large-model
escalation, projecting up to 78\% cost reduction while staying
within 2\,pp of an all-large-model baseline.  When combined with the
Visual Confused Deputy guardrail~\cite{vcd2026}, high-risk actions
escalate directly to the strongest model.

\paragraph{Adaptation regime depends on workload.}
For \emph{chat}, per-turn user feedback detected by classifiers like
our Feedback Detector~\cite{feedbackdet2026} provides a fast,
low-overhead signal to adjust routing online.  For \emph{agents},
task-level outcomes observed only after multi-step sessions call for
RL-based optimization~\cite{routerr1_2025} that reasons over
sequential routing decisions.


% -----------------------------------------------------------------------------
\subsection{Agent Loops, Memory, and Fleet-Scale Orchestration}
\label{sec:router:agent-loops}
% -----------------------------------------------------------------------------

Recent agent-serving work argues that inference for tool-using agents
should be treated as a \emph{data-systems} problem, not a sequence of
independent chat completions.  Helium~\cite{helium2026} reframes
agentic LLM serving around workflow structure, reuse, and operator
placement; Sutradhara~\cite{sutradhara2026} co-designs orchestrator
and engine for tool-heavy inference; CONCUR~\cite{concur2026}
applies congestion-aware batching to high-throughput agentic
inference; AgServe~\cite{agserve2025} shows that
\emph{session-aware} scheduling transcends static cost--quality
trade-offs; Continuum~\cite{continuum2025} couples multi-turn
scheduling with KV-cache time-to-live so that paused tool calls do
not silently destroy locality.  These systems align with
experience-driven routing (EvoRoute~\cite{evoroute2026}) and
budget-aware sequential routing (BAAR~\cite{baar2026}), but they are
typically evaluated inside a single stack.

The semantic router complements them by exporting a \emph{fleet-wide}
control surface: the same hooks that implement signal-driven model
routing can (i)~reshape tool and instruction payloads
(ITR-style~\cite{itr2026}, evaluated on that paper's controlled
benchmark), (ii)~attach memory-tier decisions (full
context vs.\ summary vs.\ retrieval-first prefill informed by
MemCost~\cite{memcost2026} and provider prompt-caching studies on
long-horizon agents~\cite{dontbreakcache2026}), and (iii)~emit scheduling hints
(admission deferral, batch affinity) that reduce queueing-induced
retries---a common hidden driver of extra agent rounds.  xRouter's
RL-trained orchestration~\cite{xrouter2025} is evidence that the
router layer can learn cost-aware policies when rewards are defined
over full trajectories; combining that with gateway-native stacks
such as OpenClaw~\cite{openclaw2026}---where real users drive
high-variance multi-turn traffic through a single control
plane---multiplies the diversity of trajectories available for
learning without requiring every framework to share one memory or
tool SDK.


